\documentclass{article}
\usepackage{geometry}
\usepackage{hyperref}
\usepackage{amsmath}
\usepackage{listings}
\geometry{a4paper, margin=1in}

\title{Pedagogical Logic and Workflows in AlgebraQuest}
\author{}
\date{\today}

\begin{document}

\maketitle

\tableofcontents

\section{Introduction}
AlgebraQuest is an adaptive intelligent tutoring system designed to enhance learning outcomes in Algebra Foundations. This report documents the pedagogical logic and workflows implemented in the system, focusing on the adaptive engine, error analysis, learner modeling, and assessment design. The system leverages data-driven insights to personalize learning experiences, ensuring that each learner receives tailored support based on their unique needs and progress.

\section{Adaptive Engine}
The adaptive engine is the core of the system, responsible for tailoring the learning experience to individual needs. It consists of several modules that work together to analyze learner behavior, identify areas of improvement, and provide targeted interventions.

\subsection{Adaptation Engine}
The adaptation engine (\texttt{adaptationEngine.js}) converts rule engine outputs and behavioral metrics into actionable decisions for the user interface. It generates a decision object with fields such as \texttt{nextAction}, \texttt{urgency}, \texttt{message}, and \texttt{recommendation}. Supported actions include:
\begin{itemize}
    \item \textbf{Micro Lesson}: Display a targeted micro-lesson to address specific misconceptions.
    \item \textbf{Increase Difficulty}: Serve a harder question to challenge the learner and promote growth.
    \item \textbf{Retry Without Hints}: Disable hint buttons for the next question to encourage independent problem-solving.
    \item \textbf{Simplify Explanation}: Provide step-by-step feedback to clarify complex concepts.
    \item \textbf{Guided Solution}: Show a fully worked example to demonstrate the correct approach.
    \item \textbf{Continue}: Proceed with normal progression when no intervention is needed.
\end{itemize}
The adaptation engine uses metrics such as confidence scores, engagement levels, and misconception scores to determine the most appropriate action. These metrics are updated dynamically based on learner interactions.

\subsection{Error Analyzer}
The error analyzer (\texttt{errorAnalyzer.js}) identifies and classifies error patterns to detect misconceptions. It uses thresholds to determine persistent and dominant errors, ensuring accurate diagnosis. Misconceptions are defined as errors that:
\begin{itemize}
    \item Occur at least three times.
    \item Account for at least 40\% of all errors.
\end{itemize}
The error analyzer tracks the following data:
\begin{itemize}
    \item \textbf{Error Frequency}: The number of times a specific error occurs.
    \item \textbf{Error Types}: Categories of errors, such as arithmetic mistakes or conceptual misunderstandings.
    \item \textbf{Error Trends}: Patterns in errors over time, indicating areas of persistent difficulty.
\end{itemize}
This data is used to refine the learner model and inform the adaptation engine.

\subsection{Learner Model}
The learner model (\texttt{learnerModel.js}) maintains a behavioral profile for each concept, tracking attempts, errors, time taken, and hint usage. It provides insights into learner behavior, such as trends and engagement levels. The following data points are tracked:
\begin{itemize}
    \item \textbf{Attempts}: The number of attempts made on each question.
    \item \textbf{Errors}: The types and frequency of errors encountered.
    \item \textbf{Time Taken}: The time spent on each question, indicating engagement and difficulty.
    \item \textbf{Hint Usage}: The number of hints requested, reflecting the learner's confidence and independence.
    \item \textbf{Concept Mastery}: A confidence score for each concept, updated based on performance.
\end{itemize}
The learner model enables the system to adapt to individual needs, providing personalized support and challenges.

\subsection{Rule Engine}
The rule engine (\texttt{ruleEngine.js}) evaluates priority-ordered rules to map conditions to pedagogical actions. Each rule is self-contained, with a condition, action, priority, and message. The context object passed to rules includes metrics like confidence score, engagement level, and misconception score. For example:
\begin{itemize}
    \item \textbf{Condition}: If the misconception score exceeds 0.7.
    \item \textbf{Action}: Display a micro-lesson.
    \item \textbf{Priority}: High.
    \item \textbf{Message}: "You seem to be struggling with this concept. Let's review it together."
\end{itemize}
The rule engine ensures that interventions are timely and relevant, addressing the most pressing needs first.

\section{Assessment Bank}
The assessment bank (\texttt{assessmentBank.js}) contains questions for each knowledge node, with correct answers, distractors, error tags, and pedagogical explanations. Each question is designed to assess specific concepts and provide meaningful feedback. For example:
\begin{itemize}
    \item \textbf{Question}: Which of the following is a variable in the expression $14 - 3y$?
    \item \textbf{Correct Answer}: $y$
    \item \textbf{Distractors}: $14$, $3$, $-3y$.
    \item \textbf{Error Tags}: Coefficient confusion, constant confusion, expression confusion.
    \item \textbf{Explanations}: Detailed feedback for each option, explaining why it is correct or incorrect.
\end{itemize}
The assessment bank is structured to support adaptive learning, with questions linked to specific knowledge nodes and misconceptions.

\section{Conclusion}
The pedagogical workflows in AlgebraQuest are designed to provide a personalized and effective learning experience. By combining adaptive decision-making, error analysis, learner modeling, and well-structured assessments, the system aims to address individual learning needs and improve outcomes. The integration of these components ensures that learners receive targeted support, enabling them to overcome challenges and achieve mastery in Algebra Foundations.

\end{document}